\section{The fixed-point setup}
\label{sec:setup}

Every pool operation triggers a sync, and every sync must publish one
number: the senior share rate $y$. \Cref{sec:overview:dawn-to-dusk}
showed why this is a fixed point, not an oracle read. This section makes
it precise: one equation in one unknown, then two facts about one
function.

\subsection{What we must prove}
\label{sec:setup:roadmap}

Write the senior NAV the sync must produce as $x$. We will see that a
self-consistent $x$ is exactly a zero of the \emph{error function}
\[
  E(x) \;=\; F\!\left(P_A(x / N)\right) - x,
\]
where $P_A$ values the junior collateral at a candidate rate and $F$ is
the accountant's sync map. The on-chain solver finds that zero by
integer bisection, and bisection needs exactly two facts:

\begin{description}
  \item[(Monotonicity)] $E$ never rises: each step holds or falls. So the
    values where $E \geq 0$ form an initial segment, the crossing is
    \emph{unique}, and bisection knows which way to step.
  \item[(Bracketing)] $E$ is non-negative at the bottom of the search range
    and negative at the top. This makes the crossing \emph{exist} and
    brackets it.
\end{description}

The rest of this section defines the three objects $F$, $P_A$, and $E$.
\Cref{sec:existence} proves the two facts. \Cref{sec:solver} runs the
solver. The argument rests on nothing else.

\subsection{The sync map $F$}
\label{sec:setup:F}

On each sync, the Dusk kernel calls the shared Royco accountant, whose
sync map we denote $F$. Given the
current raw NAVs of both tranches, $F$ runs the attribution and
settlement steps of \cref{sec:overview:waterfall} against the stored
checkpoint and returns the effective NAVs:
\[
  F\colon \bigl(\STraw,\ \JTraw;\ \text{checkpoint}\bigr) \longmapsto (\STeff,\ \JTeff).
\]
Two properties of $F$ matter downstream (\cref{sec:overview:waterfall}).
It conserves NAV ($\STraw + \JTraw = \STeff + \JTeff$), and it is
non-decreasing and 1-Lipschitz in each input (\cref{lem:F}).
Conservation holds by construction: every stored checkpoint is itself a
prior output of $F$.

In a Dusk market the senior raw NAV is exogenous to the loop: it is the
value of the senior tranche's own assets, marked independently of the
pool. Only the junior raw NAV depends on the senior rate, because the
junior's collateral is the BPT. The senior side therefore drops out, and
for one sync, with $\STraw$ and the checkpoint fixed, we read $F$ as a
function of a single variable:
\[
  F\colon \JTraw \longmapsto \STeff.
\]
$F$ returns $\STeff$ \emph{gross} of protocol fees: a fee accrues as a
separate NAV claim settled outside $F$, never netted out of $\STeff$. So
$y = \STeff/N$ is consistent only when $N$ counts the same shares the
gross NAV backs, including any minted for fees.

\subsection{The conservative valuation $P_A$}
\label{sec:setup:PA}

The junior tranche owns a fraction $f$ of a
Gyro E-CLP pool whose two legs are the senior share and a quote asset.
$P_A(y)$ is the value of that position when the senior share trades at
rate $y$, marked conservatively so that no manipulation and no rounding
can ever credit the junior tranche more than it could withdraw.

\paragraph{The frozen curve.}
The senior leg is priced by the rate the kernel publishes. At the start of a solve the
kernel reads the cached rate $y_0$ once and computes the pool invariant
$L$ from the rate-scaled balances, rounded down in favour of the senior
tranche. $L$ is then held fixed while the candidate rate $y$ varies.
Freezing $L$ is what makes $P_A$ a clean function of $y$ alone, which the
monotonicity proof relies on. The implementation details of the scaling
and rounding are collected in \cref{sec:solver:obligations}.

\paragraph{The fair point.}
Given $L$, the conservative value at candidate rate $y$ is the pool's
value at its \emph{fair point}: the point on the frozen curve whose
\emph{marginal price} (its local slope, the spot rate for an
infinitesimal trade) equals $y$. Write this value $V(y)$ and write
\[
  D(y) \;=\; V'(y),
\]
the pool's \emph{demand} for senior shares at price $y$, that is, the
senior reserve the fair point holds. Two geometric facts about the
E-CLP, proved in \cref{sec:existence:mono}, are all we need: $V$ is
non-decreasing, with $V'(y) = D(y) \geq 0$ (a higher senior price raises
the value of the senior shares the pool holds), and the demand $D(y)$ is
finite at every rate the curve can express.

\paragraph{The clamp.}
Far below the operating rate the curve's demand $D(y)$ grows without
bound, but the pool can hold at most the total supply $N$ (even if every
senior holder sold into it). The true demand cannot exceed $N$, so we
cap it there:
\[
  \Dc(y) \;=\; \min\bigl(D(y),\ N\bigr),
  \qquad
  \widehat{V}(y) \;=\; \int_0^y \Dc(u)\,du.
\]
By definition $\widehat{V}(0) = 0$, and $V(0) \geq 0$ since a reserve
is non-negative. With $\Dc \leq D = V'$ this gives
$0 \leq \widehat{V} \leq V$ everywhere: the cap only lowers the mark and
preserves conservatism. At any healthy
rate the live demand is far below $N$ and the cap is inactive. Its
purpose is to bound demand by the supply, which \cref{sec:existence}
turns into unconditional monotonicity.

\paragraph{The valuation.}
The junior tranche's pro-rata claim, rounded down in favour of the
senior tranche, is the \emph{conservative BPT valuation}:
\begin{equation}
  P_A(y) \;=\; \bigl\lfloor\, f \cdot \widehat{V}(y) \,\bigr\rfloor.
  \label{eq:PA}
\end{equation}
By the geometric facts and the clamp, $P_A$ is non-decreasing in $y$ and
Lipschitz with constant $f N$ (\cref{sec:existence:mono}). Within a
single solve $V$ depends only on $L$ and $y$, not on the live pool
composition. A fee-free swap conserves the pool's true invariant, and
the $L$ the solver consumes is a round-down underestimate of it
recomputed from live balances, so $L \leq L_{\mathrm{true}}$ for every
reachable composition. A swap (or a donation, which Balancer V3 ignores
since it accounts internally) can therefore move the mark only within
the round-down band, and never above the true fair value. The mark is
manipulation-bounded, not immovable.

\subsection{The fixed-point equation}
\label{sec:setup:circularity}

Chaining the definitions, $\JTraw = P_A(\STeff/N)$ and
$\STeff = F(\JTraw)$, so the senior NAV must reproduce itself:
\begin{equation}
  x \;=\; F\!\left(P_A\!\left(x / N\right)\right).
  \label{eq:fixed-point}
\end{equation}
A solution is a zero of the error function $E$
(\cref{sec:setup:roadmap}), and the published rate is
$y^\star = x^\star / N$. \Cref{sec:existence} proves $E$ has exactly one
zero (Monotonicity) inside a range whose endpoints straddle it
(Bracketing).
